\documentclass[10pt,a4paper]{article}

\usepackage[a4paper,top=16mm,bottom=17mm,left=17mm,right=17mm]{geometry}
\usepackage{fontspec}
\usepackage{xcolor}
\usepackage{tabularx}
\usepackage{array}
\usepackage{enumitem}
\usepackage[most]{tcolorbox}
\usepackage{fancyhdr}
\usepackage{lastpage}
\usepackage{titlesec}
\usepackage{microtype}

\IfFontExistsTF{Arial}{
  \setmainfont{Arial}
  \setsansfont{Arial}
}{
  \setmainfont{TeX Gyre Heros}
  \setsansfont{TeX Gyre Heros}
}

\definecolor{AIEPRed}{HTML}{D62027}
\definecolor{AIEPNavy}{HTML}{082B5C}
\definecolor{AIEPInk}{HTML}{182B3A}
\definecolor{AIEPMuted}{HTML}{5F6B7A}
\definecolor{AIEPLine}{HTML}{D8DEE6}
\definecolor{AIEPSoft}{HTML}{F5F2EC}
\definecolor{AIEPBlueSoft}{HTML}{E9EEF4}
\definecolor{AIEPCyan}{HTML}{35B7C6}
\definecolor{AIEPGold}{HTML}{E0BC5A}
\definecolor{AIEPGreen}{HTML}{2E8B57}

\pagestyle{fancy}
\fancyhf{}
\renewcommand{\headrulewidth}{0pt}
\fancyfoot[L]{\footnotesize\textcolor{AIEPMuted}{GeoGreen Escolar · Mentoría 2 · Planificación docente}}
\fancyfoot[R]{\footnotesize\textcolor{AIEPMuted}{\thepage/\pageref{LastPage}}}

\setlength{\parindent}{0pt}
\setlength{\parskip}{4pt}
\setlist[itemize]{leftmargin=*,topsep=2pt,itemsep=2pt,parsep=0pt}
\setlist[enumerate]{leftmargin=*,topsep=2pt,itemsep=2pt,parsep=0pt}
\renewcommand{\arraystretch}{1.16}

\titleformat{\section}{\Large\bfseries\color{AIEPNavy}}{}{0pt}{}
\titleformat{\subsection}{\normalsize\bfseries\color{AIEPNavy}}{}{0pt}{}
\titlespacing*{\section}{0pt}{2pt}{6pt}
\titlespacing*{\subsection}{0pt}{6pt}{3pt}

\newcolumntype{Y}{>{\raggedright\arraybackslash}X}
\newcolumntype{C}[1]{>{\centering\arraybackslash}p{#1}}

\newtcolorbox{docbox}[2][]{
  colback=white,
  colframe=AIEPLine,
  boxrule=0.5pt,
  arc=1.2mm,
  left=3mm,right=3mm,top=2.5mm,bottom=2.5mm,
  before skip=4pt,after skip=4pt,
  title={\scriptsize\bfseries\MakeUppercase{#2}},
  coltitle=AIEPMuted,
  fonttitle=\sffamily,
  attach title to upper={\par\vspace{1mm}},
  #1
}

\newtcolorbox{accentbox}[2][]{
  colback=AIEPSoft,
  colframe=AIEPRed,
  boxrule=0.7pt,
  arc=1.2mm,
  left=3mm,right=3mm,top=2.5mm,bottom=2.5mm,
  before skip=5pt,after skip=5pt,
  title={\scriptsize\bfseries\MakeUppercase{#2}},
  coltitle=AIEPNavy,
  fonttitle=\sffamily,
  attach title to upper={\par\vspace{1mm}},
  #1
}

\newtcolorbox{navybox}[2][]{
  colback=AIEPNavy,
  colframe=AIEPNavy,
  coltext=white,
  boxrule=0pt,
  arc=1.2mm,
  left=3.5mm,right=3.5mm,top=3mm,bottom=3mm,
  before skip=5pt,after skip=5pt,
  title={\scriptsize\bfseries\MakeUppercase{#2}},
  coltitle=AIEPGold,
  fonttitle=\sffamily,
  attach title to upper={\par\vspace{1mm}},
  #1
}

\newcommand{\eyebrow}[1]{%
  {\small\bfseries\textcolor{AIEPRed}{\MakeUppercase{#1}}}\par
  \vspace{7mm}
  {\color{AIEPRed}\rule{18mm}{0.9pt}}\par
  \vspace{9mm}
}

\newcommand{\criterion}[3]{%
  \begin{tcolorbox}[colback=#3!7,colframe=#3,boxrule=0.5pt,arc=1mm,left=2mm,right=2mm,top=1.6mm,bottom=1.6mm,before skip=2pt,after skip=2pt]
  {\bfseries\textcolor{AIEPNavy}{#1}}\hfill{\small #2}
  \end{tcolorbox}
}

\begin{document}

\eyebrow{GeoGreen Escolar · AIEP Osorno}

{\Huge\bfseries\color{AIEPNavy}Mentoría 2\par}
\vspace{2mm}
{\LARGE\bfseries\color{AIEPInk}Solución propuesta y recursos definidos\par}
\vspace{4mm}
{\large\color{AIEPMuted}Planificación docente para Ingeniería, Energía y Tecnología (I/E/T)\par}

\vspace{9mm}
\begin{tabularx}{\textwidth}{@{}>{\bfseries\color{AIEPNavy}}p{40mm}Y@{}}
Programa & GeoGreen Escolar Osorno \\
Fecha & Lunes 7 de septiembre de 2026 \\
Duración & 60 minutos por bloque \\
Participantes & 60 estudiantes, distribuidos en dos bloques de 30 \\
Organización & 5 equipos de 6 estudiantes por bloque \\
Área responsable & Ingeniería, Energía y Tecnología (I/E/T) \\
Apoyo técnico & Electricidad y Electrónica \\
Responsable & [Nombre del/de la docente] \\
\end{tabularx}

\vspace{7mm}
\begin{accentbox}{Sentido de la mentoría}
La sesión transforma un problema validado y una idea tecnológica inicial en un producto mínimo viable definido. Cada equipo debe decidir qué observará, con qué criterio actuará, cómo responderá, dónde mostrará la información y qué recursos necesita para producir evidencia antes de la Mentoría 3.
\end{accentbox}

\begin{navybox}{Entregable obligatorio de salida}
Cada equipo finaliza con una \textbf{solución funcionalmente descrita}; variable, sensor o fuente de datos y umbral definidos; \textbf{componentes, materiales y capa de software} identificados; condiciones de seguridad reconocidas; y un \textbf{plan de trabajo verificable} con responsables, fechas y evidencia comprometida.
\end{navybox}

\newpage
\section{Propósito, objetivos y condiciones de ejecución}

\begin{tabularx}{\textwidth}{@{}Y@{\hspace{5mm}}Y@{}}
\begin{docbox}[height=67mm]{Objetivo general}
Al finalizar la sesión, cada equipo podrá traducir su problema validado en un producto mínimo viable que mida una variable real, decida mediante una regla justificable, responda de forma perceptible y muestre información a un destinatario; además, podrá acotar los recursos y distribuir el trabajo necesario para producir evidencia antes del siguiente hito.
\end{docbox}
&
\begin{docbox}[height=67mm]{Objetivos específicos}
\begin{itemize}
  \item Describir la solución como observar, decidir, responder y mostrar.
  \item Seleccionar variable, sensor o fuente de datos y umbral.
  \item Definir recursos físicos, software, seguridad y alcance.
  \item Distribuir tareas entre las seis responsabilidades.
  \item Comprometer una evidencia verificable para la Mentoría 3.
\end{itemize}
\end{docbox}
\end{tabularx}

\subsection{Productos obligatorios de entrada}
Cada equipo dispone sobre la mesa de:
\begin{enumerate}
  \item el problema validado y el contexto afectado de la Mentoría 1;
  \item la ficha de idea tecnológica inicial del Taller 3;
  \item el próximo paso verificable comprometido al cierre de la Mentoría 1.
\end{enumerate}

Si falta un producto, se registra como brecha. La sesión trabaja con los antecedentes disponibles, pero no reemplaza la trazabilidad mediante una explicación improvisada.

\subsection{Condiciones que deben preservarse}
\begin{tabularx}{\textwidth}{@{}Y@{\hspace{4mm}}Y@{}}
\begin{docbox}[height=53mm]{Lógica del programa}
\begin{itemize}
  \item Cinco equipos de seis estudiantes por bloque.
  \item Se mantienen las seis responsabilidades definidas.
  \item El problema orienta la tecnología, no al revés.
  \item La salida escrita es obligatoria.
\end{itemize}
\end{docbox}
&
\begin{docbox}[height=53mm]{Continuidad}
\begin{itemize}
  \item La mentoría orienta, contrasta y devuelve decisiones.
  \item El equipo construye, prueba y registra entre sesiones.
  \item La Mentoría 3 revisa evidencia, no intenciones.
\end{itemize}
\end{docbox}
\end{tabularx}

\begin{accentbox}{Criterio de alcance}
Se prioriza una variable principal, una regla comprensible y una prueba mínima. Una solución acotada, verificable y explicable ofrece una base más sólida que varias funciones incompletas.
\end{accentbox}

\newpage
\section{Enfoque metodológico y preparación}

\begin{navybox}{Rol docente}
El/la docente ayuda a convertir intenciones en decisiones técnicas, exige fundamento para variables y umbrales, contrasta disponibilidad y seguridad, y devuelve al equipo la responsabilidad de elegir. No diseña la solución completa, no asigna componentes sin relación con el problema y no valida un montaje que no ha sido revisado físicamente.
\end{navybox}

\begin{tabularx}{\textwidth}{@{}Y@{\hspace{5mm}}Y@{}}
\begin{docbox}[height=80mm]{Principios de facilitación}
\begin{itemize}
  \item \textbf{Función antes que pieza:} definir qué debe hacer el sistema.
  \item \textbf{Variable antes que sensor:} seleccionar el componente desde la necesidad.
  \item \textbf{Regla con fundamento:} declarar de dónde proviene el umbral.
  \item \textbf{Prueba mínima:} comprobar primero una lectura que reacciona.
  \item \textbf{Seguridad presencial:} revisar el montaje real antes de energizar.
  \item \textbf{Tarea verificable:} responsable, fecha y forma de comprobación.
\end{itemize}
\end{docbox}
&
\begin{docbox}[height=80mm]{Preparación antes de la sesión}
\begin{itemize}
  \item Revisar la presentación y esta planificación.
  \item Imprimir una ficha de solución y recursos por equipo.
  \item Solicitar los tres productos de entrada.
  \item Disponer el banco acotado de sensores y salidas.
  \item Tener disponible el manual del kit de sensores.
  \item Comprobar proyector, computador y recursos visuales.
  \item Definir un punto de revisión eléctrica antes del USB.
\end{itemize}
\end{docbox}
\end{tabularx}

\subsection{Dos velocidades de trabajo}
\begin{tabularx}{\textwidth}{@{}Y@{\hspace{5mm}}Y@{}}
\begin{docbox}[height=49mm]{Velocidad alta}
Investigar, comparar componentes, interpretar documentación, calcular, planificar y construir una primera versión del software pueden acelerarse con apoyo de herramientas digitales y agentes.
\end{docbox}
&
\begin{docbox}[height=49mm]{Velocidad baja}
Voltaje, polaridad, tierra común, resistencias, conexiones y montaje se comprueban físicamente. Ante cualquier duda, el circuito permanece sin energía.
\end{docbox}
\end{tabularx}

\begin{accentbox}{Uso supervisado de agentes}
El equipo entrega contexto, solicita una tarea acotada, contrasta la respuesta y toma la decisión final. Un resultado que no puede explicar o comprobar todavía no forma parte de su proyecto.
\end{accentbox}

\newpage
\section{Plan minuto a minuto}

\small
\begin{tabularx}{\textwidth}{@{}C{17mm}p{35mm}Y p{40mm}@{}}
\textbf{\color{AIEPNavy}Tiempo} & \textbf{\color{AIEPNavy}Momento} & \textbf{\color{AIEPNavy}Acción docente} & \textbf{\color{AIEPNavy}Evidencia} \\
\hline
0--10 & \textbf{Bloque 1}\newline Del problema al producto mínimo & Recuperar los productos de entrada; orientar la cadena observar, decidir, responder y mostrar; precisar el mínimo que ya ofrece utilidad. & Descripción funcional inicial y destinatario de la información. \\
\hline
10--25 & \textbf{Bloque 2}\newline Variable, sensor y umbral & Pedir unidad o estados observables; contrastar sensores o fuentes de datos; exigir el origen del umbral; hacer leer la regla completa. & Variable, componente, umbral fundamentado y regla escrita. \\
\hline
25--45 & \textbf{Bloque 3}\newline Recursos, software y seguridad & Ordenar recursos disponibles, por conseguir y representados; definir qué mostrará el software; revisar seguridad, tiempo y alcance. & Lista de recursos, nivel de software, alcance y alternativa. \\
\hline
45--57 & \textbf{Bloque 4}\newline Trabajo y evidencia & Traducir las seis responsabilidades en tareas verificables; elegir una rama de profundización y precisar qué se mostrará el 21 de septiembre. & Plan con personas, fechas, comprobación y evidencia comprometida. \\
\hline
57--60 & \textbf{Cierre}\newline Resguardar decisiones & Verificar la ficha completa y pedir la declaración final del equipo. & Ficha de solución y recursos conservada por el equipo. \\
\end{tabularx}
\normalsize

\vspace{6mm}
\begin{accentbox}{Administración del tiempo}
La sesión no se utiliza para construir el proyecto completo. Cada bloque deja una decisión registrada y reduce incertidumbre. El avance material, las pruebas, el software y el registro continúan con trabajo del equipo entre sesiones.
\end{accentbox}

\subsection{Rondas docentes focalizadas}
\begin{tabularx}{\textwidth}{@{}p{37mm}Y@{}}
\textbf{Ronda 1 · coherencia} & ¿Qué parte del problema resuelve cada función propuesta? \\
\textbf{Ronda 2 · medición} & ¿Qué variable cambia, en qué unidad y de dónde proviene el umbral? \\
\textbf{Ronda 3 · viabilidad} & ¿Qué existe, qué falta, qué se representará y qué requiere revisión? \\
\textbf{Ronda 4 · compromiso} & ¿Quién hará qué, para cuándo y cómo podrá comprobarse? \\
\end{tabularx}

\begin{navybox}{Regla temporal}
El trabajo crítico se organiza entre el 7 y el 11 de septiembre. La semana del 14 no contempla actividades del programa y la Mentoría 3 se realiza el 21 de septiembre.
\end{navybox}

\newpage
\section{Guía práctica de facilitación}

\subsection{Bloque 1 · de la idea a una función comprobable}
\begin{tabularx}{\textwidth}{@{}Y@{\hspace{5mm}}Y@{}}
\begin{docbox}[height=63mm]{Preguntas clave}
\begin{itemize}
  \item ¿Qué observa el sistema?
  \item ¿Con qué condición decide?
  \item ¿Qué respuesta puede percibirse?
  \item ¿Dónde se muestra la información?
  \item ¿Quién la utiliza y para qué?
\end{itemize}
\end{docbox}
&
\begin{docbox}[height=63mm]{Señales de avance}
\begin{itemize}
  \item La solución puede explicarse sin enumerar piezas.
  \item El mínimo mide, decide, responde y muestra.
  \item Existe un destinatario reconocible.
  \item La propuesta conserva relación con el problema validado.
\end{itemize}
\end{docbox}
\end{tabularx}

\subsection{Bloque 2 · variable, componente y regla}
La secuencia de orientación es:
\begin{enumerate}
  \item \textbf{Precisar:} identificar una variable principal, unidad o dos estados distinguibles.
  \item \textbf{Comparar:} revisar componentes desde la variable y las condiciones de uso.
  \item \textbf{Fundamentar:} obtener el umbral mediante medición, cálculo, fuente o acuerdo provisional declarado.
  \item \textbf{Formular:} escribir y leer la regla sin explicaciones adicionales.
\end{enumerate}

\begin{navybox}{Estructura de la regla}
Cuando \rule{27mm}{0.35pt} pase de \rule{21mm}{0.35pt}, entonces \rule{28mm}{0.35pt}, para que \rule{24mm}{0.35pt} pueda \rule{25mm}{0.35pt}.
\end{navybox}

\subsection{Bloque 3 · recursos y seguridad}
\begin{tabularx}{\textwidth}{@{}p{39mm}Y@{}}
\textbf{Ya disponible} & Componentes, herramientas, conocimientos, permisos y espacios confirmados. \\
\textbf{Debe conseguirse} & Recurso, responsable, fecha y alternativa si no llega. \\
\textbf{Se representará} & Simulación, maqueta, diagrama o registro cuando construir no sea viable. \\
\textbf{Software} & Qué muestra, a quién, con qué datos y dónde debe funcionar. \\
\end{tabularx}

\begin{accentbox}{Revisión antes de energizar}
Comprobar voltaje, polaridad, tierra común, resistencias y montaje. El relé y cualquier intervención sobre 220 V requieren supervisión docente. Un montaje dudoso no se conecta.
\end{accentbox}

\newpage
\section{Plan de trabajo, evidencias y continuidad}

\subsection{Bloque 4 · convertir responsabilidades en tareas}
\begin{tabularx}{\textwidth}{@{}p{36mm}Y@{}}
\textbf{Coordinación} & Calendario, turnos, acuerdos y seguimiento. \\
\textbf{Investigación} & Fuentes para el sensor, el umbral y el problema. \\
\textbf{Diseño} & Boceto, montaje, ubicación y experiencia de uso. \\
\textbf{Tecnología} & Diagrama, conexiones, lógica y primera versión del software. \\
\textbf{Pruebas y evidencia} & Plan de prueba, mediciones y registro. \\
\textbf{Comunicación} & Explicación comprensible, fotografías y materiales ordenados. \\
\end{tabularx}

\begin{navybox}{Formato de compromiso}
\textbf{Qué} se realizará · \textbf{quién} lo asume · \textbf{para cuándo} · \textbf{cómo se comprueba}. Cada responsabilidad deja al menos una tarea verificable.
\end{navybox}

\subsection{Verificadores del entregable}
\criterion{Coherencia}{La solución responde al problema y se describe por funciones.}{AIEPRed}
\criterion{Medición}{Variable, unidad o estados, sensor y umbral están definidos.}{AIEPCyan}
\criterion{Viabilidad}{Recursos, software, seguridad, alcance y alternativa son explícitos.}{AIEPGold}
\criterion{Organización}{Las tareas tienen persona, fecha y forma de comprobación.}{AIEPGreen}
\criterion{Evidencia}{El equipo sabe qué mostrará y qué aprendizaje explicará.}{AIEPRed}

\subsection{Evidencia aceptable para la Mentoría 3}
Prototipo parcial funcionando, simulación, diagrama verificado, maqueta instrumentada, registro de mediciones o panel digital. Debe responder con claridad: \textbf{qué probamos} y \textbf{qué aprendimos}.

\begin{tabularx}{\textwidth}{@{}Y@{\hspace{5mm}}Y@{}}
\begin{docbox}[height=53mm]{El equipo conserva}
\begin{itemize}
  \item regla completa del producto mínimo viable;
  \item lista de recursos y condiciones de seguridad;
  \item nivel de software y rama de profundización;
  \item plan de trabajo y alternativa;
  \item evidencia comprometida.
\end{itemize}
\end{docbox}
&
\begin{docbox}[height=53mm]{Cierre verbal}
\textbf{Nuestro producto mínimo viable va a:}
\vspace{3mm}
\hrule
\vspace{5mm}
\textbf{Antes del 21 de septiembre demostraremos que:}
\vspace{3mm}
\hrule
\end{docbox}
\end{tabularx}

\begin{accentbox}{Continuidad obligatoria}
La ficha de solución y recursos se convierte en el plan vigente del equipo. Entre sesiones se ejecutan las tareas, se realizan pruebas y se conserva evidencia. La Mentoría 3 revisa lo construido, probado y registrado, incluyendo errores y ajustes.
\end{accentbox}

\end{document}
